%----------------------------------------------------------------------------------------
%   USEFUL COMMANDS
%----------------------------------------------------------------------------------------

\newcommand{\dipartimento}{Dipartimento di Matematica ``Tullio Levi-Civita''}

%----------------------------------------------------------------------------------------
% 	USER DATA
%----------------------------------------------------------------------------------------

% Data di approvazione del piano da parte del tutor interno; nel formato GG Mese AAAA
% compilare inserendo al posto di GG 2 cifre per il giorno, e al posto di 
% AAAA 4 cifre per l'anno
\newcommand{\dataApprovazione}{Data}

% Dati dello Studente
\newcommand{\nomeStudente}{Antonio}
\newcommand{\cognomeStudente}{Sandu}
\newcommand{\matricolaStudente}{2113194}
\newcommand{\emailStudente}{antonio.sandu@studenti.unipd.it}
\newcommand{\telStudente}{+ 39 348 65 62 943}

% Dati del Tutor Aziendale
\newcommand{\nomeTutorAziendale}{Fabio}
\newcommand{\cognomeTutorAziendale}{Pallaro}
\newcommand{\emailTutorAziendale}{f.pallaro@synclab.it}
\newcommand{\telTutorAziendale}{+ 39 333 13 68 500}
\newcommand{\ruoloTutorAziendale}{Manager}

% Dati dell'Azienda
\newcommand{\ragioneSocAzienda}{SyncLab S.p.A}
\newcommand{\indirizzoAzienda}{Galleria Spagna, 28, Padova (PD)}
\newcommand{\sitoAzienda}{https://synclab.it/}
\newcommand{\emailAzienda}{info@synclab.it}
\newcommand{\partitaIVAAzienda}{P.IVA 07952560634}

% Dati del Tutor Interno (Docente)
\newcommand{\titoloTutorInterno}{Prof.}
\newcommand{\nomeTutorInterno}{Luigi}
\newcommand{\cognomeTutorInterno}{De Giovanni}

\newcommand{\prospettoSettimanale}{
     % Personalizzare indicando in lista, i vari task settimana per settimana
     % sostituire a XX il totale ore della settimana
    \begin{itemize}
        \item \textbf{Prima Settimana - Discussione requisiti e Ripasso teorico (32 ore)}
        \begin{itemize}
            \item incontro con persone coinvolte nel progetto per discutere 
            i requisiti e le richieste relativamente al progetto da sviluppare;
            \item ripasso Java Standard Edition e tool di sviluppo (IDE, StopLight);
            \item ripasso principi della buona programmazione (SOLID, CleanCode);
            \item studio teorico dell'architettura a microservizi: passaggio da monolite 
            ad architetture a microservizi con pro e contro (Api Gateway, Service
            Discovery e Service Registry, Circuit Breaker e Saga).
        \end{itemize}
        \item \textbf{Seconda Settimana - Autoformazione Backend (40 ore)} 
        \begin{itemize}
            \item studio Spring Core/Spring Boot;
            \item studio ORM, in particolare il framework Spring Data JPA;
            \item studio servizi REST e framework spring Data REST;
            \item realizzazione mini prototipo di progetto SpringBoot con endpoint REST.
        \end{itemize}
        \item \textbf{Terza Settimana - Autoformazione Frontend (40 ore)} 
        \begin{itemize}
            \item ripasso linguaggio Javascript;
            \item studio del linguaggio TypeScript;
            \item studio piattaforma NodeJS e AngularCLI;
            \item studio framework Angular;
            \item implementazione di un mini prototipo di una maschera in Angular con 
            5 campi generici di input con chiamata al back-end tramite un service dedicato.
        \end{itemize}
        \item \textbf{Quarta Settimana - Progettazione architetturale (40 ore)} 
        \begin{itemize}
            \item progettazione architetturale ed entità del progetto da realizzare;
            \item sviluppo e test API REST del back-end dei servizi per le CRUD di gestione utenti e partite.
        \end{itemize}
        \item \textbf{Quinta Settimana - Sviluppo e test delle maschere di interfaccia (40 ore)} 
        \begin{itemize}
            \item sviluppo e test delle maschere di interfaccia di frontend per 
            le CRUD che richiamano il back-end precedentemente realizzato.
        \end{itemize}
        \item \textbf{Sesta Settimana - Codifica (40 ore)} 
        \begin{itemize}
            \item sviluppo della parte main del gioco (gestione della partita) sia per la
            parte di back end che front end.
        \end{itemize}
        \item \textbf{Settima Settimana - Conclusione della codifica (40 ore)} 
        \begin{itemize}
            \item conclusione sviluppi e loro integrazione.
        \end{itemize}
        \item \textbf{Ottava Settimana - Conclusione (40 ore)} 
        \begin{itemize}
            \item stesura del documento di Specifica Tecnica;
            \item ultimo incontro di validazione degli artefatti con il tutor aziendale.
        \end{itemize}
    \end{itemize}
}

% Indicare il totale complessivo (deve essere compreso tra le 300 e le 320 ore)
\newcommand{\totaleOre}{312}

\newcommand{\obiettiviObbligatori}{
	 \item \underline{\textit{O01}}: acquisizione di competenze pratiche nell'utilizzo dei framework Angular e Spring attraverso attività di autoformazione e sviluppo di prototipi applicativi;
	 \item \underline{\textit{O02}}: definizione dell'architettura e dei requisiti funzionali della web app attraverso attività di analisi dei requisiti e progettazione;
	 \item \underline{\textit{O03}}: implementazione del backend della web app mediante il framework Spring e realizzazione di API REST dedicate alla gestione delle funzionalità applicative;
	 \item \underline{\textit{O04}}: implementazione del frontend della web app mediante il framework Angular, realizzando un'interfaccia responsive per l'interazione con le funzionalità del sistema;
     \item \underline{\textit{O05}}: integrazione delle componenti frontend e backend realizzando un MVP (Minimum Viable Product) funzionante conforme ai requisiti definiti durante le fasi iniziali del progetto;
     \item \underline{\textit{O06}}: redazione di un documento di Specifica Tecnica che descriva architettura, tecnologie utilizzate e principali scelte progettuali adottate durante lo sviluppo.
}

\newcommand{\obiettiviDesiderabili}{
	 \item \underline{\textit{D01}}: containerizzazione di almeno una parte tra backend e frontend del progetto;
	 \item \underline{\textit{D02}}: apprendimento e applicazione di strumenti e modalità operative adottate nel contesto aziendale per la gestione e lo sviluppo del progetto.
}

\newcommand{\obiettiviFacoltativi}{
	 \item \underline{\textit{F01}}: containerizzazione di tutte le parti del progetto.
}